\begin{tikzpicture}
  \node[panel, inner sep=8pt] (card) {%
    \begin{minipage}{0.95\textwidth}
    \setlength{\parskip}{0pt}
    {\footnotesize\sffamily\bfseries\color{inkmute} NOTATION: THE NINE QUANTITIES THAT TRAVEL THROUGH BOTH PARTS}\par\smallskip
    \footnotesize
    \begin{tabular}{@{}l@{\quad}p{0.50\textwidth}@{\quad}p{0.17\textwidth}@{}}
    \toprule
    symbol & reading & fixed at \\
    \midrule
    \(\Time\) & machine steps to halt: \(\Time_M(x)\) on one input, and
      \(\Time_D(n)=\max_{x,y,a,b}\Time_{D,(n,x,y,a,b)}\) for a decider &
      \S1.1; Definition~\ref{def:paper-decider} \\
    \(|M|,|t|\) & description size in bits, with
      \(|t|=|\operatorname{ser}(t)|\) for a term & \S1.4;
      Definition~\ref{def:l-serialization} \\
    \(f,F\) & remaining fuel, and a fuel budget supplied to one bounded run &
      \S8.3 \\
    \(\kappa_p\) & the \emph{charge} of one primitive step; every microstep
      spends one fuel unit & \S8.3 \\
    \(\operatorname{level}\) & the CL nesting depth;
      \(\operatorname{level}(V)=\operatorname{level}(S)\), and it is an upper
      bound, never a minimum &
      Definitions~\ref{def:paper-lambda} and~\ref{def:cl-func-doc} \\
    \(\lambda\) & the resource bound: \(|V|\leq\lambda\) and
      \(\Time_S(n),\Time_D(n)\leq n^\lambda\); also \(\Compress\)'s second
      argument & Definition~\ref{def:paper-lambda} \\
    \(Q\) & a bound on the public question lengths \(|x|,|y|\); the
      answer-reduction contract sets \(Q(n)=(\lambda n)^\mu\) &
      Theorem~\ref{thm:l-succinct-sat} \\
    \(R\) & \(R=2^{r_0}\), the succinct index range of the three tableau
      blocks of a decoupled 5SAT instance & \S11.3 \\
    \(k(n)\) & \(k(n)=(\lambda n)^{(1+c_3')\tau}\), the number of copies
      \(\Repeat\) runs & Figure~\ref{fig:parameter-card} \\
    \bottomrule
    \end{tabular}\par\smallskip
    {\scriptsize\sffamily\color{inkmute}
     The remaining constants are collected where they are first needed:
     \(\mu,\gamma,\tau,C_0\) and the level \(9\) in
     Figure~\ref{fig:parameter-card}; \(h(d,u)=3+|d|+|\operatorname{enc}(u)|\)
     and \(c_Y=3\) in \S8.4 and \S10.2; \(C_L\) and the exponent \(8\) in
     Theorem~\ref{thm:l-to-tm}; \(A\) in
     Theorem~\ref{thm:lambda-preservation}; \(q=2^k\), \(m\), \(m'\), \(d\)
     and \(\degF\) in \S11.5 and \S12.  Two names are deliberately kept apart:
     \(\tau\) is always the repetition constant, and the transition-window
     function of Lemma~\ref{lem:transition-window} is \(\omega\).}
    \end{minipage}};
\end{tikzpicture}
